\chapter{Analytic Control and Noncollapsing}

This chapter establishes the analytic toolkit for controlled Ricci flows with
surgery.  Reduced geometry, local curvature estimates, compactness, and
noncollapsing provide the inputs used to analyze singularity models and to
continue the flow in the next chapter.

\section{Surgery Flows and Fixed Controls}

\begin{definition}[Ricci flow with surgery]
\label{def:ricci-flow-with-surgery}
\dcref{MT:ch14:14.8}
\notready
A \emph{Ricci flow with surgery} is a pair $(\mathcal M,G)$ with the
following structure.  The surgery spacetime $\mathcal M$ has a continuous
time function $\mathbf t\colon\mathcal M\to I$, a time vector field $\chi$
with $\chi(\mathbf t)=1$, and the ordinary, one-sided, and surgery charts
which distinguish smooth, exposed, and singular points.  Every slice
$M_t=\mathbf t^{-1}(t)$ is a smooth manifold, and the horizontal bundle is
$\mathcal H=\ker(d\mathbf t)$.  The tensor $G$ is a smooth metric on
$\mathcal H$ and satisfies
\[
 \mathcal L_\chi G=-2\operatorname{Ric}_G.
\]
At an exposed point or an endpoint of $I$, the Lie derivative is one-sided.
The flow is three-dimensional when its slices are three-manifolds.
\end{definition}

\begin{definition}[Noncollapsing for a surgery flow]
\label{def:surgery-flow-noncollapsing}
\dcref{MT:ch14:14.14,MT:ch14:14.15}
\uses{def:ricci-flow-with-surgery}
\notready
Let $(\mathcal M,G)$ be a three-dimensional Ricci flow with surgery.  The
backward parabolic neighborhood $P(x,t,\rho,-\rho^2)$ exists when every
point of $B_{g(t)}(x,\rho)$ has a backward $\chi$-flow line of length
$\rho^2$, and these flow lines give a product embedding of the ball over
$[t-\rho^2,t]$.  For $\kappa,r_0>0$, the flow is
\emph{$\kappa$-noncollapsed on scales at most $r_0$} if
\[
 P(x,t,\rho,-\rho^2)\text{ exists},\qquad
 |\operatorname{Rm}|\le \rho^{-2}\text{ on }P
 \quad\Longrightarrow\quad
 \operatorname{Vol}_{g(t)}B(x,t,\rho)\ge\kappa\rho^3
\]
whenever $0<\rho\le r_0$.
\end{definition}

\begin{definition}[Strong canonical-neighborhood control]
\label{def:strong-canonical-neighborhood-control}
\dcref{MT:ch9:9.76,MT:ch14:14.18}
\uses{def:ricci-flow-with-surgery,def:epsilon-neck-and-scale,def:canonical-cap}
\notready
Fix $C<\infty$, $0<\epsilon<1/2$, and $r>0$.  A three-dimensional Ricci
flow with surgery satisfies the \emph{strong
$(C,\epsilon)$-canonical-neighborhood assumption with parameter $r$} if
every point $x$ with $R(x)\ge r^{-2}$ has one of the following
neighborhoods at the curvature scale $R(x)^{-1/2}$:
\begin{enumerate}
\item a strong $\epsilon$-neck centered at $x$, whose parabolic rescaling by
      $R(x)$ is $\epsilon$-close in
      $C^{\lfloor\epsilon^{-1}\rfloor}$ to the shrinking round cylinder on
      normalized times $(-1,0]$;
\item a $(C,\epsilon)$-cap whose core contains $x$, with $C$-comparable
      scalar curvatures and scale-normalized diameter and volume bounded by
      $C$, together with a uniform core volume lower bound and the
      scale-invariant derivative bounds for $R$;
\item a compact positively curved $C$-component containing $x$, with
      least sectional curvature at least $C^{-1}\sup R$ and diameter between
      $C^{-1}(\sup R)^{-1/2}$ and $C(\inf R)^{-1/2}$;
\item an $\epsilon$-round component containing $x$: a compact connected
      component $(M,g)$ for which there are $R>0$, a compact
      constant-curvature-one manifold $(Z,g_0)$, and a diffeomorphism
      $\varphi:Z\to M$ such that $\varphi^*(Rg)$ is $\epsilon$-close to $g_0$
      in the $C^{\lfloor\epsilon^{-1}\rfloor}$ topology.
\end{enumerate}
Only the first alternative is a spacetime neighborhood.  The other
alternatives lie in the slice $M_{\mathbf t(x)}$.
\end{definition}

\begin{definition}[Controlled surgery parameters]
\label{def:controlled-surgery-parameters}
\dcref{MT:ch15:15.5,MT:ch15:15.6,MT:ch15:15.7}
\uses{def:strong-canonical-neighborhood-control,thm:local-neck-surgery}
\notready
Fix $\epsilon>0$ sufficiently small and choose $C$ larger than the uniform
canonical-neighborhood constants for the standard solution and for
three-dimensional $\kappa$-solutions.  Put $t_0=2^{-5}$ and
$T_i=2^it_0$.  The control data are positive nonincreasing sequences
\[
 \mathbf r=(r_i)_{i\ge0},\qquad
 \mathbf K=(\kappa_i)_{i\ge0},\qquad
 \boldsymbol\delta=(\delta_i)_{i\ge0},
 \qquad r_0=\epsilon,
\]
and positive nonincreasing piecewise constant functions $r(t)$ and
$\kappa(t)$ with values $r_{i+1}$ and $\kappa_{i+1}$ on
$[T_i,T_{i+1})$.  A positive nonincreasing surgery-control function
$\bar\delta(t)$ is required to satisfy
$\bar\delta(t)\le\delta_{i+1}$ on that interval.  Define
\[
 \rho(t)=\bar\delta(t)r(t),
 \qquad
 h(t)=h(\rho(t),\bar\delta(t))
 \le \bar\delta(t)^2r(t),
\]
where $h$ is the neck scale at which the local surgery construction applies.
At a surgery time $t$, every cutting sphere is the central sphere of a strong
$\bar\delta(t)$-neck of scalar curvature $h(t)^{-2}$.
\end{definition}

\begin{definition}[Controlled Ricci flow with surgery]
\label{def:controlled-ricci-flow-with-surgery}
\dcref{MT:ch15:15.9,MT:ch15:15.11}
\uses{def:ricci-flow-with-surgery,def:strong-canonical-neighborhood-control,def:controlled-surgery-parameters,def:maximal-ricci-flow-extension}
\notready
For the fixed $C$, $\epsilon$, and control functions, a
\emph{controlled Ricci flow with surgery} is a three-dimensional Ricci flow
with surgery satisfying the following conditions.
\begin{enumerate}
\item Every slice is compact and contains no embedded $\mathbb{RP}^2$ with
      trivial normal bundle.
\item Singular times form a discrete subset of the time interval.
\item Time zero is initial, and every component of $(M_0,g(0))$ is
      normalized: $|\operatorname{Rm}|\le1$ and every unit ball has at least
      half the Euclidean unit-ball volume.
\item At a singular time $t$, the exposed region $E_t$ is a finite disjoint
      union of three-balls, and
      $M_t=C_t\cup_{\Sigma_t}E_t$ with $\Sigma_t=\partial E_t$ a finite
      union of two-spheres.
\item Every component of $\Sigma_t$ is a central sphere at the prescribed
      scale in the maximal presurgery extension to $t$.
\item Immediately before $t$, every point which disappears has a strong
      $(C,\epsilon)$-canonical neighborhood.
\item Between consecutive singular times the ordinary Ricci flow is maximal.
\end{enumerate}
The flow is additionally required to have curvature pinched toward positive
and to satisfy the strong canonical-neighborhood assumption with parameter
$r(t)$.  Its stagewise noncollapsing condition is that, whenever
$0<\rho\le\epsilon$ and
$P(x,t,\rho,-\rho^2)$ exists with
$|\operatorname{Rm}|\le\rho^{-2}$, either
\[
 \operatorname{Vol}_{g(t)}B(x,t,\rho)\ge\kappa(t)\rho^3
\]
or the component of $M_t$ containing $x$ has positive sectional curvature.
\end{definition}

\begin{theorem}[Local volume and injectivity-radius control]
\label{thm:local-volume-injectivity-radius-control}
\dcref{MT:ch1:1.35,MT:ch1:1.36,MT:ch5:5.10}
\notready
Fix $n$, $K<\infty$, $r_0>0$, and $v_0>0$.  There are positive constants
\[
 v_*=v_*(n,K,r_0,v_0),
 \qquad \iota_*=\iota_*(n,K,r_0,v_0),
\]
with the following property.  If $(M^n,g)$ is complete,
$|\operatorname{Rm}|\le K$ on $B(p,r_0)$, and
$\operatorname{Vol}B(p,r_0)\ge v_0$, then, for
$0<s\le\min\{r_0,K^{-1/2}\}$,
\[
 \operatorname{Vol}B(p,s)\ge v_*s^n,
 \qquad
 \operatorname{inj}_g(p)\ge
 \iota_*\min\{r_0,K^{-1/2}\}.
\]
The convention is that $K^{-1/2}=\infty$ when $K=0$.
\end{theorem}

\begin{proof}
Put $a=\min\{r_0,K^{-1/2}\}$ and rescale by $a^{-2}$, so the relevant
ball has radius one and sectional curvatures have absolute value at most one.
In particular $\operatorname{Ric}\ge-(n-1)$.  Bishop--Gromov comparison with
hyperbolic $n$-space gives, for $0<s\le1$,
\[
 \frac{\operatorname{Vol}B(p,s)}{V_{-1}(s)}
 \ge
 \frac{\operatorname{Vol}B(p,1)}{V_{-1}(1)}.
\]
Applying the same comparison first between radii $a$ and $r_0$ before
rescaling supplies a positive lower bound for
$\operatorname{Vol}B(p,1)$ depending only on the displayed data.  Since
$V_{-1}(s)/s^n$ has a positive minimum on $(0,1]$, this proves the asserted
volume ratio.

For the injectivity estimate, curvature at most one excludes conjugate points
on radial geodesics of length less than one.  If the injectivity radius at
$p$ tended to zero with the volume ratio fixed, Klingenberg's alternative
would therefore give a shortest geodesic loop of length twice the injectivity
radius.  Lift the unit ball to the nonsingular exponential ball in $T_pM$.
The successive lifts obtained by concatenating this loop are distinct until
their total length reaches a fixed fraction of one.  Curvature comparison
bounds from below the volume of one fixed lifted subball, while the distinct
lifts all project into $B(p,1)$ with uniformly bounded multiplicity.  Hence
the volume of $B(p,1)$ is at most a dimensional constant times the
injectivity radius, contradicting the positive volume lower bound.  This
gives a positive rescaled injectivity bound; scaling back gives $\iota_*a$.
\end{proof}

\begin{theorem}[Normalized short-time flow]
\label{thm:normalized-short-time-ricci-flow}
\dcref{MT:ch3:3.11,MT:ch4:4.11,MT:ch5:5.16,MT:ch15:15.1}
\notready
There is a universal $\kappa_0>0$ such that every normalized compact
three-manifold $(M,g_0)$ admits a unique Ricci flow $g(t)$ on
$[0,2^{-5}]$ with $g(0)=g_0$.  On this interval
\[
 |\operatorname{Rm}_{g(t)}|\le2,
 \qquad
 \operatorname{Vol}_{g(t)}B_{g(t)}(x,\rho)\ge\kappa_0\rho^3
 \quad(0<\rho\le1).
\]
Consequently this interval has no surgery time and supplies the initial
stage of a controlled Ricci flow with surgery.
\end{theorem}

\begin{proof}
\uses{def:ricci-flow-with-surgery,def:controlled-ricci-flow-with-surgery,def:maximal-ricci-flow-extension,thm:local-volume-injectivity-radius-control}
The Ricci--DeTurck equation obtained by adding the Lie derivative along the
trace of the difference between the connections of $g$ and $g_0$ is strictly
parabolic.  Short-time parabolic existence, followed by pullback under the
generating diffeomorphisms, gives a Ricci flow; the reverse gauge argument
gives uniqueness.

The curvature evolution inequality and $|\operatorname{Rm}_{g_0}|\le1$ give
$|\operatorname{Rm}_{g(t)}|\le2$ for $t\le2^{-4}$ as long as the flow
exists.  On the same interval metric and volume distortion give universal
constants $a,b>0$ such that
\[
 B_{g_0}(x,a\rho)\subset B_{g(t)}(x,\rho),
 \qquad
 d\operatorname{vol}_{g(t)}\ge b\,d\operatorname{vol}_{g_0}.
\]
Apply \ref{thm:local-volume-injectivity-radius-control} at time zero.
Normalization and $|\operatorname{Rm}_{g_0}|\le1$ give one universal lower
volume ratio $v_*>0$ on all radii at most one.  The distortion estimates therefore
yield
\[
 \operatorname{Vol}_{g(t)}B_{g(t)}(x,\rho)
 \ge b\operatorname{Vol}_{g_0}B_{g_0}(x,a\rho)
 \ge bv_*a^3\rho^3.
\]
Take $\kappa_0=bv_*a^3$.  If the maximal flow ended before
$2^{-5}$, its uniform curvature bound and the bounded-curvature extension
theorem would extend it, a contradiction.  Thus it exists through $2^{-5}$;
viewed as a surgery flow, its surgery set is empty.
\end{proof}

\section{Reduced Geometry}

\begin{definition}[Reduced length]
\label{def:reduced-length}
\dcref{MT:ch6:6.2,MT:ch6:6.45}
\uses{def:ricci-flow-with-surgery}
\notready
Fix a smooth point $(x,T)$ and write backward time as $\tau=T-t$.  For a
piecewise smooth backward path $\gamma\colon[0,\bar\tau]\to\mathcal M$
with $\gamma(0)=x$, define
\[
 \mathcal L(\gamma)=
 \int_0^{\bar\tau}\sqrt\tau
 \bigl(R(\gamma(\tau))+|\dot\gamma(\tau)|^2\bigr)\,d\tau.
\]
For a point $y$ at backward time $\bar\tau$, let
$L_x(y,\bar\tau)$ be the infimum of $\mathcal L$ over such paths and put
\[
 l_x(y,\bar\tau)=\frac{L_x(y,\bar\tau)}{2\sqrt{\bar\tau}}.
\]
These quantities are used only on domains on which every competing path
stays in the smooth locus.
\end{definition}

\begin{definition}[Reduced exponential map and injectivity domain]
\label{def:reduced-exponential-map}
\dcref{MT:ch6:6.17,MT:ch6:6.25,MT:ch6:6.26}
\uses{def:reduced-length}
\notready
An $\mathcal L$-geodesic from $(x,T)$ is a critical point of
$\mathcal L$.  Its initial vector is
\[
 v=\lim_{\tau\downarrow0}\sqrt\tau\,\dot\gamma(\tau)
 \in T_xM_T.
\]
The reduced exponential map is
$\mathcal L\exp_x^\tau(v)=\gamma_v(\tau)$ whenever this geodesic exists.
Its injectivity set $\mathcal D_x^\tau$ consists of those $v$ for which
$\gamma_v|_{[0,\tau]}$ is the unique minimizing $\mathcal L$-geodesic to its
endpoint and has no nonzero $\mathcal L$-Jacobi field vanishing at both ends.
The image $\Omega_x^\tau=\mathcal L\exp_x^\tau(\mathcal D_x^\tau)$ is the
reduced injectivity domain.
\end{definition}

\begin{definition}[Reduced volume]
\label{def:reduced-volume}
\dcref{MT:ch6:6.70,MT:ch6:6.71}
\uses{def:reduced-length,def:reduced-exponential-map}
\notready
When the reduced injectivity domain has full measure in the time slice, the
reduced volume based at $(x,T)$ is
\[
 \widetilde V_x(\tau)=
 \int_{M_{T-\tau}}(4\pi\tau)^{-3/2}e^{-l_x(y,\tau)}
 \,d\operatorname{vol}_{g(T-\tau)}(y).
\]
Equivalently, if $J(v,\tau)$ is the Jacobian of
$\mathcal L\exp_x^\tau$, then
\[
 \widetilde V_x(\tau)=
 \int_{\mathcal D_x^\tau}(4\pi\tau)^{-3/2}
 e^{-l_x(\mathcal L\exp_x^\tau(v),\tau)}J(v,\tau)\,dv.
\]
\end{definition}

\begin{lemma}[Upper support inequality for reduced length]
\label{lem:reduced-length-support-inequality}
\dcref{MT:ch6:6.4,MT:ch6:6.13,MT:ch6:6.49,MT:ch6:6.51,MT:ch7:7.8}
\notready
On a smooth bounded-curvature Ricci-flow domain, reduced length is smooth
before the reduced cut locus and satisfies, in dimension three,
\[
 \partial_\tau l+\Delta l\le \frac{3/2-l}{\tau}.
\]
At an arbitrary endpoint of a minimizing $\mathcal L$-geodesic and for every
$\eta>0$, there is a smooth upper support $l_\eta$ which agrees with $l$ at
that endpoint and satisfies
\[
 \partial_\tau l_\eta+\Delta l_\eta
 \le \frac{3/2-l_\eta}{\tau}+\eta.
\]
\end{lemma}

\begin{proof}
\uses{def:reduced-length,def:ricci-flow-with-surgery}
For a variation field $Y$ along a minimizing curve $\gamma$, differentiate
the defining integral in \ref{def:reduced-length}.  Integration by parts,
using $\partial_tg=-2\operatorname{Ric}$, gives the endpoint term and the
$\mathcal L$-geodesic equation
\[
 \nabla_X X-\frac12\nabla R+2\operatorname{Ric}(X,\cdot)^*
 +\frac{1}{2\tau}X=0.
\]
Differentiate once more.  Trace the index form over fields with orthonormal
endpoint values and compare them with
$Y_i(s)=\sqrt{s/\tau}\,E_i(s)$, where $E_i$ is parallel after the Ricci
correction.  The boundary terms combine with the first-variation identities
for $l=L/(2\sqrt\tau)$, and the traced inequality is exactly
$\partial_\tau l+\Delta l\le(3/2-l)/\tau$.

At a cut endpoint, stop a fixed minimizing curve at a time
$0<\tau_0<\tau$.  In a neighborhood of its remaining segment, the reduced
exponential map is nonsingular.  Add the fixed $\mathcal L$-length of the
initial segment to the smooth terminal value function.  The resulting
function is at least $l$, equals it at the chosen endpoint, and obeys the
same second-variation calculation with an error tending to zero as
$\tau_0\downarrow0$.  Choosing $\tau_0$ so that this error is below $\eta$
gives the asserted support.
\end{proof}

\begin{lemma}[Minimum reduced-length comparison]
\label{lem:minimum-reduced-length-comparison}
\dcref{MT:ch7:7.10,MT:ch7:7.11,MT:ch7:7.12}
\notready
Suppose that on every slice of a smooth backward-time interval the minimum
\[
 m(\tau)=\min_y l_x(y,\tau)
\]
is attained, and that the union of the minimizing sets over compact
subintervals is compact.  Then $m$ is continuous and
\[
 D^+m(\tau)\le\frac{3/2-m(\tau)}{\tau}.
\]
If $\limsup_{\tau\downarrow0}m(\tau)\le3/2$, then
$m(\tau)\le3/2$ throughout the interval.
\end{lemma}

\begin{proof}
\uses{lem:reduced-length-support-inequality,lem:forward-dini-mean-value}
Compactness of the minimizing set makes the upper supports in
\ref{lem:reduced-length-support-inequality} uniform for nearby slices.
Transport a minimizer by the time vector field for the upper estimate and a
nearby minimizer backward for the lower estimate.  The two estimates prove
continuity and, after division by the time increment and passage to the
limsup, the displayed Dini inequality.

Put $q=m-3/2$.  On every compact interval bounded away from zero,
$D^+q\le-q/\tau$.  If $q$ became positive, apply
\ref{lem:forward-dini-mean-value} to $e^{\int d\tau/\tau}q=\tau q$ between
the last zero and a later positive value.  The product has nonpositive upper
right Dini derivative but would have to increase, a contradiction.  Letting
the left endpoint decrease to zero proves $m\le3/2$.
\end{proof}

\begin{theorem}[Regularity of the reduced exponential map]
\label{thm:reduced-exponential-regularity}
\dcref{MT:ch6:6.18,MT:ch6:6.28,MT:ch6:6.56,MT:ch6:6.67}
\notready
Assume that minimizing $\mathcal L$-geodesics to an open set remain in a
smooth region on which $|\operatorname{Ric}|$ and $|\nabla R|$ are bounded.
Then $\mathcal D_x^\tau$ is open,
\[
 \mathcal L\exp_x^\tau\colon\mathcal D_x^\tau
 \longrightarrow\Omega_x^\tau
\]
is a diffeomorphism, and $l_x(\cdot,\tau)$ is smooth on
$\Omega_x^\tau$.  The complement of $\Omega_x^\tau$ in the selected open
set has measure zero.
\end{theorem}

\begin{proof}
\uses{def:reduced-length,def:reduced-exponential-map,lem:reduced-length-support-inequality}
The $\mathcal L$-geodesic equation is a smooth ODE after the change of
variable $s=\sqrt\tau$.  Its derivative with respect to the initial vector is
the $\mathcal L$-Jacobi equation.  Absence of a Jacobi field vanishing at the
two endpoints makes this derivative invertible, so the inverse-function
theorem gives local smoothness of $\mathcal L\exp_x^\tau$.  Uniqueness of the
minimizer makes the local inverses agree, proving the diffeomorphism and the
smoothness of $l$.

The curvature bounds and the first-variation formula give a local Lipschitz
bound for $l$.  An endpoint with two distinct minimizing geodesics is a point
where $l$ has two different endpoint gradients, hence is not differentiable;
these points form a null set by Rademacher's theorem.  An endpoint with a
conjugate minimizer is a critical value of the smooth reduced exponential map
and therefore forms a null set by Sard's theorem.  Every point outside the
injectivity domain is of one of these two types, which proves the last claim.
\end{proof}

\begin{lemma}[Reduced-density monotonicity]
\label{lem:reduced-density-monotonicity}
\dcref{MT:ch6:6.50,MT:ch6:6.80,MT:ch6:6.81}
\notready
Let $v\in\mathcal D_x^\tau$ and let $J(v,\tau)$ be the Jacobian of the
reduced exponential map.  Along the corresponding minimizing
$\mathcal L$-geodesic,
\[
 \frac d{d\tau}\left((4\pi\tau)^{-3/2}e^{-l}J\right)\le0.
\]
Consequently the reduced volume is nonincreasing in backward time as long as
all minimizing curves under consideration remain in one smooth domain.
\end{lemma}

\begin{proof}
\uses{def:reduced-volume,lem:reduced-length-support-inequality,thm:reduced-exponential-regularity}
Choose three $\mathcal L$-Jacobi fields whose endpoint values are an
orthonormal basis.  The logarithmic derivative of their volume is
$\partial_\tau\log J$.  The traced second-variation estimate used in
\ref{lem:reduced-length-support-inequality}, together with the endpoint
identities for $\partial_\tau l$, gives
\[
 \partial_\tau\log J-\partial_\tau l-\frac{3}{2\tau}\le0.
\]
This is the logarithmic derivative of the displayed density.  By
\ref{thm:reduced-exponential-regularity}, the injectivity domain has full
measure, so change of variables in \ref{def:reduced-volume} and monotone
convergence give the reduced-volume assertion without a cut-locus boundary
term.
\end{proof}

\begin{theorem}[Quantitative reduced-volume transfer]
\label{thm:reduced-volume-transfer}
\dcref{MT:ch6:6.71,MT:ch6:6.80,MT:ch8:8.1}
\notready
Let $x$ be the terminal point at time $T$, let $0<\bar\tau<T$, and let
$E$ be a measurable subset of the regular reduced-exponential image at
backward time $\bar\tau$.  Let
\[
 d\mu_{\bar\tau}(y)=(4\pi\bar\tau)^{-3/2}
 e^{-l_x(y,\bar\tau)}\,d\operatorname{Vol}_{g(T-\bar\tau)}(y)
 \quad (y\in E).
\]
Assume the minimizing curves defining $l_x$ remain in one compact smooth
domain.  Suppose there is a measurable endpoint map
$F:E\to M_T$ and that the push-forward $\mu_T=F_*\mu_{\bar\tau}$ is
absolutely continuous on a terminal ball, with density $\theta_T$ satisfying
 \[
 \operatorname{Vol}_{g(T-\bar\tau)}(E)\ge V>0,
 \qquad l_x(\cdot,\bar\tau)\le L_0\quad\hbox{on }E,
 \]
and
\[
 \operatorname{supp}\mu_T\subset B_{g(T)}(x,\rho),
 \qquad 0\le\theta_T\le C\rho^{-3}\quad\hbox{a.e.}
 \]
Then
\[
 \operatorname{Vol}_{g(T)}B_{g(T)}(x,\rho)
 \ge C^{-1}(4\pi\bar\tau)^{-3/2}e^{-L_0}V\,\rho^3.
\]
\end{theorem}

\begin{proof}
The lower bound on $l_x$ gives
\[
 \mu_{\bar\tau}(E)\ge(4\pi\bar\tau)^{-3/2}e^{-L_0}V.
\]
Push-forward preserves total mass, while the density and support hypotheses
give
\[
 \mu_T(M_T)=\int_{B_{g(T)}(x,\rho)}\theta_T\,d\operatorname{Vol}_{g(T)}
 \le C\rho^{-3}\operatorname{Vol}_{g(T)}B_{g(T)}(x,\rho).
\]
Combining the two inequalities proves the claim.  In applications, the
absolute-continuity and density bound are the substantive endpoint-transfer
hypotheses; they are not implied by the fact that all reduced geodesics have
the same terminal basepoint.
\end{proof}

\section{Local Curvature and Neck Estimates}

\begin{theorem}[Local Ricci-flow derivative estimates]
\label{thm:local-ricci-flow-derivative-estimates}
\dcref{MT:ch3:3.28,MT:ch3:3.33}
\notready
For every $m\ge0$ there is $C_m<\infty$ with the following property.  If a
Ricci flow has a relatively compact parabolic neighborhood
$P(x,t,r,-r^2)$ on which $|\operatorname{Rm}|\le r^{-2}$, then
\[
 |\nabla^m\operatorname{Rm}|
 \le C_m r^{-m-2}
\]
on $P(x,t,r/2,-r^2/2)$.
\end{theorem}

\begin{proof}
\uses{def:ricci-flow-with-surgery}
Parabolically rescale so that $r=1$.  Commuting covariant derivatives with
the curvature evolution equation gives
\[
 (\partial_t-\Delta)\nabla^m\operatorname{Rm}
 =\sum_{a+b=m}\nabla^a\operatorname{Rm}*\nabla^b\operatorname{Rm}.
\]
For $m=1$, apply the maximum principle to
$t|\nabla\operatorname{Rm}|^2+A|\operatorname{Rm}|^2$ after multiplying by
a spatial cutoff supported in the unit ball; choosing $A$ larger than the
universal reaction coefficients absorbs the cross terms.  This gives a
uniform first-derivative bound on the half cylinder.

Inductively, assume the estimates through order $m-1$.  Apply the same
argument to
\[
 t^m|\nabla^m\operatorname{Rm}|^2
 +\sum_{j<m}A_jt^j|\nabla^j\operatorname{Rm}|^2,
\]
using nested cutoffs.  The induction bounds every lower-order product, and
the constants $A_j$ absorb the commutator terms.  The parabolic maximum
principle gives a bound depending only on $m$.  Scaling back yields the
factor $r^{-m-2}$.
\end{proof}

\begin{theorem}[Strong maximum principle for curvature tensors]
\label{thm:strong-curvature-tensor-maximum-principle}
\dcref{MT:ch4:4.17,MT:ch4:4.18}
\uses{thm:time-slice-maximum-comparison}
\notready
Let $\mathcal T$ solve
$\partial_t\mathcal T=\Delta\mathcal T+\Psi(\mathcal T)$ on a compact
parabolic domain, and let $K$ be a fiberwise convex, parallel-invariant set.
If $\Psi$ lies in the tangent cone of $K$ along $\partial K$, and the initial
and parabolic-boundary values lie in $K$, then $\mathcal T$ remains in $K$.
If an invariant supporting functional is strictly positive at one interior
initial point, it is strictly positive at every later interior point.
\end{theorem}

\begin{proof}
Choose a nonnegative Dirichlet heat solution $h$ supported near the point of
strict positivity and initially below the supporting functional.  The scalar
comparison theorem \ref{thm:time-slice-maximum-comparison} makes $h$ positive
in the later interior.  For $\eta>0$, parallel transport the supporting
half-space inward by $\eta h$.  At a first contact of $\mathcal T$ with its
boundary, the spatial Laplacian points into the convex set, the reaction term
is tangent or inward, and the time derivative of the displaced supporting
plane points strictly outward.  This contradicts first contact.  Letting
$\eta\downarrow0$ proves preservation; keeping $\eta>0$ proves the strict
interior assertion.
\end{proof}

\begin{theorem}[Splitting at a curvature null direction]
\label{thm:curvature-null-splitting}
\dcref{MT:ch4:4.19,MT:ch4:4.20,MT:ch4:4.21}
\notready
Let a connected three-dimensional Ricci flow have nonnegative curvature
operator.  If its curvature operator has a null eigenvector at a positive
time, then the tangent Ricci-null line distribution is locally parallel and time-independent,
and every connected parabolic neighborhood on which it is defined locally
splits as a surface flow times a flat interval.  If the time slices are
complete, the universal cover splits isometrically as a Ricci flow on a
surface times a flat line.  Every quotient is obtained from this product by
a group preserving the parallel line distribution.  In
particular, if the scalar curvature vanishes at one positive-time point,
then the flow is flat on that connected spacetime neighborhood.
\end{theorem}

\begin{proof}
\uses{thm:strong-curvature-tensor-maximum-principle}
At the prescribed null point write the curvature-operator eigenvalues as
$\lambda\ge\mu\ge\nu=0$.  Extend a unit null two-form by spatial parallel
transport and by the standard time-dependent orthonormal gauge.  Its
curvature quadratic form has a spacetime minimum at that point.  The normal
component of the three-dimensional reaction term there is
$\nu^2+\lambda\mu=\lambda\mu$.  The one-sided time derivative at the final
time is nonpositive and the spatial Laplacian is nonnegative, so the
curvature evolution equation forces $\lambda\mu=0$.

If $R=0$ at a point, nonnegativity makes every curvature-operator eigenvalue
zero there.  Apply the strict assertion of
\ref{thm:strong-curvature-tensor-maximum-principle} to the scalar curvature:
positivity at any earlier point would force positivity at the prescribed
point.  Thus the curvature operator vanishes on every earlier connected
slice, and the scalar evolution propagates this flatness backward as well.

In the nonflat case $\lambda>0$, so the preceding reaction calculation gives
$\mu=0$ at the prescribed point.  Let $s=\mu+\nu$, the sum of the two
smallest curvature-operator eigenvalues.  If $s$ were positive at an earlier
interior point, the strict assertion of
\ref{thm:strong-curvature-tensor-maximum-principle} would make it positive at
the prescribed point.  Hence $s$ vanishes throughout the connected earlier
region.  Equality in the curvature reaction ODE gives spectrum
$\{\lambda,0,0\}$ there, and the scalar strong principle gives
$\lambda>0$ everywhere.

The two-dimensional kernel of the curvature operator corresponds to a
one-dimensional tangent kernel $\mathcal D$ of the Ricci tensor.  Choose a
local unit section $V$ of $\mathcal D$.  Fix $x$, $X\in T_xM$, and a curve
through $x$ tangent to $X$.  Parallel-extend $V(x)$ and an arbitrary vector
$W$ along this curve.  Nonnegative sectional curvature makes
\[
 \operatorname{Rm}(\widetilde V,\widetilde W,
                        \widetilde V,\widetilde W)\ge0,
\]
with equality at $x$; its first variation therefore vanishes.  Polarization
gives
$(\nabla_X\operatorname{Rm})(V,\cdot,\cdot,\cdot)=0$.
Differentiate
$\operatorname{Rm}(V,W_1,W_2,W_3)=0$ along three parallel fields $W_i$.
The preceding identity yields
\[
 2\operatorname{Rm}(\nabla_XV,W_1,W_2,W_3)=0.
\]
The kernel description makes $\nabla_XV$ a multiple of $V$, whereas
$|V|=1$ makes it orthogonal to $V$; hence $\nabla_XV=0$.

Thus the metric locally splits as a positively curved surface times a flat
line, and its curvature tensor is pulled back from the surface.  The
curvature evolution equation has the same product form, so
$(\partial_t\operatorname{Rm})(V,\cdot,\cdot,\cdot)=0$; differentiating the
kernel identity in time gives $\partial_tV$ in the null line, and the
unit-length condition makes it zero in the time-dependent orthonormal gauge.
The distribution is therefore parallel and time-independent.  The de Rham
theorem splits the simply connected complete cover along this line, and deck
transformations preserve it, giving the asserted global and quotient forms.
\end{proof}

\begin{theorem}[Ancient scalar-curvature monotonicity]
\label{thm:ancient-scalar-curvature-monotonicity}
\dcref{MT:ch4:4.37,MT:ch4:4.39}
\notready
Let $(M,g(t))$, $-\infty<t\le0$, be a complete Ricci flow with bounded
nonnegative curvature operator on each time slice.  Then
\[
 \partial_tR(x,t)\ge0
\]
at every spacetime point.
\end{theorem}

\begin{proof}
Hamilton's differential Harnack inequality, applied on $(T_0,0]$ with zero
vector field, gives
\[
 \partial_tR+\frac{R}{t-T_0}\ge0.
\]
For fixed $(x,t)$ the curvature is finite and nonnegative.  Letting
$T_0\to-\infty$ makes the second term tend to zero and proves the assertion.
\end{proof}

\begin{theorem}[Weak superharmonicity of Busemann functions]
\label{thm:busemann-function-superharmonicity}
\dcref{MT:ch2:2.3}
\notready
Let $(M^n,g)$ be complete with nonnegative Ricci curvature, and let $B$ be
the Busemann function of a ray.  Then $\Delta B\le0$ weakly: for every
nonnegative compactly supported smooth function $\varphi$,
\[
 \int_M B\,\Delta\varphi\,d\mu\le0.
\]
\end{theorem}

\begin{proof}
For the ray $\gamma$, put
$B_T(x)=d(x,\gamma(T))-T$.  Laplacian comparison away from the cut locus,
interpreted by upper barriers, gives
\[
 \int_M B_T\Delta\varphi\,d\mu
 \le \int_M\frac{n-1}{d(x,\gamma(T))}\varphi\,d\mu.
\]
If the support of $\varphi$ lies in the $D$-ball about $\gamma(0)$, the
denominator is at least $T-D$.  The right side therefore tends to zero.
The functions $B_T$ converge locally uniformly to $B$ by monotonicity and
the triangle inequality, so passage to the limit proves the weak inequality.
\end{proof}

\begin{lemma}[Long minimizing geodesics align with a neck axis]
\label{lem:epsilon-neck-geodesic-axis-alignment}
\dcref{MT:ch19:19.4}
\uses{def:epsilon-neck-and-scale}
\notready
For every $\alpha>0$ there is $\epsilon_0(\alpha)>0$ such that, in a
normalized $\epsilon$-neck with $\epsilon\le\epsilon_0$, every minimizing
geodesic of length greater than $\epsilon^{-1}/100$ which crosses in the
positive axial direction has tangent within $\alpha$ of the axial unit
vector at every point of the geodesic contained in the neck.
\end{lemma}

\begin{proof}
Suppose not and take $\epsilon_j\downarrow0$.  Pull the metrics and geodesics
back to the standard cylinders.  On every fixed compact subcylinder the
metrics converge in $C^2$ to the product metric and the geodesics converge
to a minimizing product geodesic.  Its axial displacement tends to infinity,
whereas its spherical displacement is at most the diameter of the round
two-sphere.  Recenter at any chosen point of misalignment.  On the resulting
full- or half-cylinder limit, the conserved product momenta force its
spherical speed to zero and its positive axial speed to one.  Smooth
convergence of the geodesic equation gives the same estimate for the
approximating tangent at the recentered point, a contradiction.
\end{proof}

\begin{lemma}[Distance comparison from outside a thin neck]
\label{lem:external-point-neck-distance-comparison}
\dcref{MT:ch19:19.9}
\uses{def:epsilon-neck-and-scale,lem:epsilon-neck-geodesic-axis-alignment}
\notready
For every $0<\alpha<1$ there is $\epsilon_2(\alpha)>0$ such that the
following holds.  Let $N$ be a separating $\epsilon$-neck,
$\epsilon\le\epsilon_2$, let $z$ lie outside its middle two-thirds on the
negative side.  Assume every minimizing geodesic from $z$ to a point of the
positive half of $N$ enters $N$ through the negative end and remains in $N$
until that endpoint.  Let $p$ lie in its middle half, and let $q$ lie anywhere
later on the same axial product line in $N$, with $s(q)>s(p)$.  Then
\[
 (1-\alpha)r_N(s(q)-s(p))
 \le d(z,q)-d(z,p)
 \le(1+\alpha)r_N(s(q)-s(p)).
\]
\end{lemma}

\begin{proof}
The upper bound is the triangle inequality and the cylindrical metric
estimate along the axial segment.  For the lower bound, take a minimizing
geodesic from $z$ to $q$.  The added hypothesis forces it to enter from the
negative end and to remain in the neck.  Once it reaches the fiber through
$p$, the portion already traversed inside the neck has length greater than
$r_N\epsilon^{-1}/100$ after shrinking $\epsilon$.
By \ref{lem:epsilon-neck-geodesic-axis-alignment}, its tangent has axial
component at least $1-\alpha/2$.  Cut the geodesic at the fiber through $p$
and compare the remaining segment with axial displacement.  The neck-metric
error is at most another $\alpha/2$ of that displacement, giving the lower
bound.
\end{proof}

\begin{lemma}[Metric spheres cross a thin neck in a fiber sphere]
\label{lem:metric-sphere-crosses-neck-as-fiber}
\dcref{MT:ch19:19.10}
\notready
For every $\alpha>0$ there is $\epsilon_1(\alpha)>0$ with the following
property.  Let $N$ be an $\epsilon$-neck, $\epsilon\le\epsilon_1$, whose
central sphere separates the ambient connected manifold.  If $z$ lies
outside the middle two-thirds of $N$ and $p$ lies in its middle sixth, then
$\partial B(z,d(z,p))\cap N$ is a topological two-sphere in the middle
quarter of $N$, and axial projection maps it homeomorphically to a neck
fiber.
\end{lemma}

\begin{proof}
\uses{lem:external-point-neck-distance-comparison}
Scale the neck to scalar curvature one.  On each axial product line,
\ref{lem:external-point-neck-distance-comparison} makes distance from $z$
strictly increasing through the middle half, with a uniform two-sided slope.
The comparison at the boundary fibers, together with the uniformly bounded
diameter of a neck fiber, puts $d(z,p)$ strictly between the two boundary
values on every axial line.  Thus each axial line meets the level set exactly
once in the middle quarter.  Applying the comparison with $q$ anywhere in
the neck shows that outside the middle quarter the axial displacement, of
order $\epsilon^{-1}$, dominates the fiber-diameter error, so there are no
additional points of the level set elsewhere in $N$.
The distance function is continuous, and the uniform strict-monotonicity
estimate implies directly that the unique axial coordinate depends
continuously on the $S^2$ coordinate: otherwise two convergent fibers and
the slope bound would give distinct limiting roots on one fiber.  The level
set is therefore a continuous graph over $S^2$, and axial projection is a
homeomorphism.  The same boundary comparisons place the graph in the middle
quarter.  No differentiability of the distance function at its cut locus is
used.
\end{proof}

\section{Compactness and Noncollapsing}

\begin{theorem}[Pointed Riemannian compactness]
\label{thm:pointed-riemannian-compactness}
\dcref{MT:ch5:5.6,MT:ch5:5.9,MT:ch5:5.10}
\notready
Let $(M_j,g_j,x_j)$ be complete pointed Riemannian manifolds.  Suppose that
on every fixed based ball all covariant derivatives of curvature are
uniformly bounded and that for some $r_0,v_0>0$,
\[
 \operatorname{Vol}_{g_j}B(x_j,r_0)\ge v_0.
\]
Then a subsequence converges smoothly on compact subsets to a complete
pointed Riemannian manifold.
\end{theorem}

\begin{proof}
\uses{thm:local-volume-injectivity-radius-control}
The base-ball curvature and volume bounds give a positive lower bound for
the injectivity radius at $x_j$ by
\ref{thm:local-volume-injectivity-radius-control}.  Bishop--Gromov comparison
and the curvature bounds propagate the same type of volume lower bound to
centers in each fixed based ball, so the same theorem gives a uniform
injectivity-radius bound there.

Exponential charts of a fixed radius therefore cover every such ball with a
uniform number of charts.  In harmonic refinements of these charts, the
curvature-derivative hypotheses give uniform $C^k$ bounds for the metric
coefficients.  Arzela--Ascoli and a diagonal choice in the radius and $k$
give compatible limiting charts and smooth transition maps.  They assemble
to the pointed smooth limit.  Every finite-length limiting geodesic remains
inside one ball of the exhaustion and is the limit of extendible geodesics
in $M_j$; hence it extends.  Hopf--Rinow gives completeness.
\end{proof}

\begin{theorem}[Pointed Ricci-flow compactness]
\label{thm:hamilton-compactness-generalized-flows}
\dcref{MT:ch5:5.14,MT:ch5:5.15}
\notready
Let $(\mathcal M_j,G_j,x_j)$ be based generalized Ricci flows with base time
zero, defined on a common open time interval $I$.  Suppose that for every
$A<\infty$ and compact $J\Subset I$ containing zero, all sufficiently large
$j$ admit compatible embeddings
\[
 B_{g_j(0)}(x_j,A)\times J\longrightarrow\mathcal M_j
\]
with compact zero-time balls and a uniform curvature bound on the images.
Assume that for some $r_0,\kappa>0$,
\[
 \operatorname{Vol}_{g_j(0)}B(x_j,r_0)\ge\kappa r_0^3.
\]
After passing to a subsequence, the flows converge smoothly on compact
spacetime subsets.  The limit is a pointed Ricci flow
\[
 (M_\infty,g_\infty(t),x_\infty)
\]
on $I$, and its time-zero slice is complete.  If
$\operatorname{Ric}_{g_\infty(t)}\ge0$ for $t\le0$, then every nonpositive
time slice of the limit is complete.
\end{theorem}

\begin{proof}
\uses{def:ricci-flow-with-surgery,def:surgery-flow-noncollapsing,thm:local-ricci-flow-derivative-estimates,thm:pointed-riemannian-compactness}
The curvature bound and the Ricci-flow equation give uniform bounds for all
covariant derivatives of curvature on smaller compact cylinders by
\ref{thm:local-ricci-flow-derivative-estimates}.  The base-ball volume lower
bound and curvature bounds put the zero-time slices under
\ref{thm:pointed-riemannian-compactness}, producing a complete smooth
zero-time limit and compatible exhaustion charts.

Exhaust each zero-time slice by the compact balls and exhaust $I$ by compact
intervals.  Repeating the construction and diagonalizing produces compatible
limit charts; their transition maps agree because they are limits of the
original embeddings.  The limiting metrics satisfy the Ricci-flow equation
by smooth convergence, and pointed Riemannian compactness already gives
completeness at time zero.

If the limiting Ricci curvature is nonnegative and $t\le0$, then
$g_\infty(t)\ge g_\infty(0)$ as quadratic forms.  A finite-length
$g_\infty(t)$-curve escaping every compact set would therefore have finite
$g_\infty(0)$-length, contradicting completeness of the time-zero slice.
Thus every nonpositive slice is complete under the additional curvature
hypothesis.  Local based-ball bounds alone make no completeness assertion at
other times.
\end{proof}

\begin{lemma}[Nonflat terminal cones are excluded]
\label{lem:nonflat-terminal-cone-exclusion}
\dcref{MT:ch4:4.22}
\notready
Let a three-dimensional Ricci flow have nonnegative sectional curvature on a
connected spacetime neighborhood.  If at a positive time a nonempty open
subset of a slice is isometric to an open subset of a metric cone, then that
spacetime neighborhood is flat.
\end{lemma}

\begin{proof}
\uses{thm:curvature-null-splitting}
A nonflat three-dimensional cone has curvature operator of rank one away
from its vertex: its null two-plane contains the radial direction, while its
nonzero eigenvalue varies quadratically under radial dilation.  By
\ref{thm:curvature-null-splitting}, rank-one degeneracy at a positive time
makes the null distribution parallel and splits the flow locally as a
surface times a line.  The nonzero curvature
eigenvalue is then constant along the parallel line, contrary to its radial
variation on a nonflat cone.  Thus the cone is flat.  The scalar-zero clause
of \ref{thm:curvature-null-splitting} propagates flatness throughout the
connected spacetime neighborhood.
\end{proof}

\begin{lemma}[Short reduced paths avoid surgery caps]
\label{lem:short-reduced-paths-avoid-surgery-caps}
\dcref{MT:ch16:16.5,MT:ch16:16.8,MT:ch16:16.13,MT:ch16:16.15}
\notready
Fix a controlled stage, a terminal point $(x,T)$ with a controlled
parabolic ball of radius at least $r_{i+1}$, and $L_0<\infty$.  There is a
positive surgery-control bound, depending only on the stage data and $L_0$,
such that every backward path from $x$ with
\[
 \mathcal L_+(\gamma)=\int\sqrt\tau
   \bigl(\max(R,0)+|\dot\gamma|^2\bigr)\,d\tau<L_0
\]
misses a fixed open neighborhood of every surgery cap and every exposed
point in $[T_{i-1},T]$.
\end{lemma}

\begin{proof}
\uses{def:controlled-ricci-flow-with-surgery,def:reduced-length,thm:hamilton-compactness-generalized-flows,def:standard-surgery-cap-model,lem:local-surgery-standard-cap-convergence}
For $0<a<b$, Cauchy--Schwarz gives
\[
 \int_a^b\sqrt\tau|\dot\gamma|^2d\tau
 \ge\frac{\sqrt a\,d(\gamma(a),\gamma(b))^2}{b-a}.
\]
After rescaling a newly inserted cap by its surgery scale $h$, the initial
metrics converge by \ref{lem:local-surgery-standard-cap-convergence} to the
standard cap.  The curvature and derivative bounds, together with
\ref{thm:hamilton-compactness-generalized-flows}, extend this convergence on
each fixed ball through every normalized time interval $[0,\theta]$ with
$\theta<1$.

Choose a standard annulus so wide that crossing it makes the right side of
the displayed estimate greater than $2L_0$.  A path which does not cross the
annulus but remains in the inner cap until normalized time $\theta$ has
\[
 \int_a^b\sqrt\tau R\,d\tau
 \ge c\sqrt a\int_a^b
 \frac{h^{-2}\,dt}{1-(t-\bar t)h^{-2}},
\]
which exceeds $2L_0$ when $\theta$ is sufficiently close to one.  Smooth
convergence preserves a lower bound $L_0$ in either alternative once the
surgery control is small.  The terminal controlled ball excludes entry from
the other end for very small backward time.  Hence a path of the stated
length cannot enter the chosen cap neighborhoods.  The exposed points are
the limiting boundaries of the same neighborhoods and are avoided as well.
\end{proof}

\begin{proposition}[Compact smooth domain for reduced minimizers]
\label{prop:compact-reduced-minimizing-domain}
\dcref{MT:ch16:16.4,MT:ch16:16.21,MT:ch16:16.24,MT:ch16:16.25}
\notready
Under the hypotheses of \ref{lem:short-reduced-paths-avoid-surgery-caps},
there is an open subset $U$ of the smooth spacetime over
$[T_{i-1},T)$ such that:
\begin{enumerate}
\item every point of $U$ is joined to $x$ by a minimizing smooth
      $\mathcal L$-geodesic;
\item every slice $U_t$ is nonempty and
      $\min_{U_t}\mathcal L\le3\sqrt{T-t}$;
\item over every compact time interval, the slice minimizers form a compact
      subset of the smooth spacetime.
\end{enumerate}
\end{proposition}

\begin{proof}
\uses{lem:short-reduced-paths-avoid-surgery-caps,lem:minimum-reduced-length-comparison}
The scalar lower bound $R\ge-6$ gives
\[
 \mathcal L_+(\gamma)
 \le\mathcal L(\gamma)+4T_{i+1}^{3/2}.
\]
Remove the cap neighborhoods supplied by
\ref{lem:short-reduced-paths-avoid-surgery-caps}.  There are only finitely
many singular times on the compact interval already constructed, and every
slice is compact, so the complement $Y(L_0)$ is compact and smooth.  It
contains every path whose $\mathcal L$-length is below the chosen bound.

A minimizing sequence in $Y(L_0)$ has a uniformly bounded weighted kinetic
energy.  Away from backward time zero this gives equicontinuity by
Cauchy--Schwarz; inside the terminal controlled ball the curvature and metric
distortion estimates give the same conclusion up to zero.  Arzela--Ascoli
gives uniform convergence, weak $W^{1,2}$ compactness gives lower
semicontinuity of the kinetic term, and smooth convergence in $Y(L_0)$ gives
convergence of the scalar term.  The limit is therefore a smooth minimizing
$\mathcal L$-geodesic.

Let $U$ consist of endpoints joined to $x$ by paths below one fixed bound
larger than $8\sqrt{T_{i+1}}(1+T_{i+1})$.  The direct method gives the first
claim, and the same compact set with varying endpoint time gives the third.
Near backward time zero, the controlled terminal ball has minimum reduced
length at most $3/2$.  Apply
\ref{lem:minimum-reduced-length-comparison}; if the interval on which this
bound holds ended early, compactness of its minimizers and a short vertical
extension would continue it.  Thus the bound holds on every slice and gives
the second claim.
\end{proof}

\begin{lemma}[Reduced-volume noncollapsing transfer]
\label{lem:reduced-volume-noncollapsing-transfer}
\dcref{MT:ch16:16.1,MT:ch16:16.27,MT:ch6:6.80}
\notready
Assume the stage-$i$ noncollapsing bound on earlier controlled slices and the
absolute-continuity/density hypotheses of \ref{thm:reduced-volume-transfer}
for the selected minimizing families.  Under the hypotheses of
\ref{prop:compact-reduced-minimizing-domain}, there is
$\kappa_{\mathcal L}>0$, depending only on the stage data, such that every
terminal parabolic ball of radius
$\rho\in[r_{i+1},\epsilon]$ and curvature at most $\rho^{-2}$ satisfies
\[
 \operatorname{Vol}B(x,T,\rho)\ge\kappa_{\mathcal L}\rho^3,
\]
unless the component containing $x$ has positive sectional curvature.
\end{lemma}

\begin{proof}
\uses{def:controlled-ricci-flow-with-surgery,def:strong-canonical-neighborhood-control,prop:compact-reduced-minimizing-domain,thm:reduced-volume-transfer}
Choose a minimizing $\mathcal L$-geodesic from $x$ to time $T_{i-1}$ with
length at most $3\sqrt{T-T_{i-1}}$.  If its scalar curvature were at least
$r_i^{-2}$ throughout the middle time interval, integration of
$\sqrt\tau R$, together with $R\ge-6$ on the two end intervals, would give
length at least $4\sqrt{T-T_{i-1}}$.  Hence the curve contains a point $y$
with $R(y)<r_i^{-2}$.

The canonical derivative estimates and curvature pinching give a backward
parabolic ball of a fixed radius $r'>0$ about $y$ with
$|\operatorname{Rm}|\le(r')^{-2}$.  A surgery cap cannot meet this ball:
its curvature is bounded below by a fixed multiple of
$\bar\delta^{-4}r_i^{-2}$, while the ball has curvature at most
$2r_i^{-2}$.  The induction hypothesis therefore gives either a
positive-curvature component or
$\operatorname{Vol}B(y,r')\ge\kappa_i(r')^3$.

In the second case, join the minimizing curve to all points of this earlier
ball by short spatial segments.  Metric distortion gives a common upper
bound $L_0$ for their reduced lengths.  The minimizing curves supplied by
\ref{prop:compact-reduced-minimizing-domain} lie in one compact smooth
domain.  Applying \ref{thm:reduced-volume-transfer} with
$V=\kappa_i(r')^3$ yields the claimed terminal lower bound.  In the first
case, preservation of positive sectional curvature carries the stated
exception to the component of $x$.
\end{proof}

\begin{theorem}[Conditional noncollapsing extension]
\label{thm:conditional-noncollapsing-extension}
\dcref{MT:ch16:16.1}
\uses{thm:reduced-volume-transfer}
\notready
Assume the endpoint-transfer hypotheses of \ref{thm:reduced-volume-transfer}
are available uniformly for the minimizing families used at scales
$[r_{i+1},\epsilon]$.  Fix $i\ge0$ and positive nonincreasing finite sequences
\[
 \delta_0\ge\cdots\ge\delta_i,
 \qquad \epsilon=r_0\ge\cdots\ge r_i,
 \qquad \kappa_0\ge\cdots\ge\kappa_i.
\]
There is $0<\kappa_{i+1}\le\kappa_i$ such that, for every
$0<r_{i+1}\le r_i$, there is
$0<\delta(r_{i+1})\le\delta_i$ with the following property.  If a
controlled Ricci flow with surgery on $[0,T)$,
$T\in(T_i,T_{i+1}]$, satisfies the stage-$i$ induction hypotheses before
$T_i$, has $\bar\delta\le\delta(r_{i+1})$ on the new stage, and satisfies
the strong canonical-neighborhood assumption with parameter $r_{i+1}$,
then every parabolic ball of radius at most $\epsilon$ satisfying the
curvature hypothesis in
\ref{def:surgery-flow-noncollapsing} has volume at least
$\kappa_{i+1}\rho^3$, unless its center lies in a positively curved component
of the time slice.  The constant is independent of the particular flow.
\end{theorem}

\begin{proof}
\uses{def:surgery-flow-noncollapsing,def:strong-canonical-neighborhood-control,def:controlled-ricci-flow-with-surgery,thm:hamilton-compactness-generalized-flows,lem:reduced-length-support-inequality,lem:reduced-volume-noncollapsing-transfer}
At scales below $r_{i+1}$, inspect the canonical neighborhood of the center.
A neck has a uniform cylindrical volume ratio, and a cap contains at a
comparable curvature scale a ball whose volume is bounded below by a constant
depending only on $C$ and $\epsilon$.  A controlled or round component gives
the positive-curvature alternative.  Thus the nonexceptional cases give a
uniform number $\kappa_{\mathrm{can}}>0$ at these scales.  The compactness
used to make this number uniform is
\ref{thm:hamilton-compactness-generalized-flows}.

For $r_{i+1}\le\rho\le\epsilon$, decrease
$\delta(r_{i+1})$ so that the surgery-cap barrier and the compact minimizing
domain apply.  The support inequality
\ref{lem:reduced-length-support-inequality} controls the minimizing reduced
length, and \ref{lem:reduced-volume-noncollapsing-transfer} gives the fixed
lower volume ratio $\kappa_{\mathcal L}$.  Taking
\[
 \kappa_{i+1}=min\{\kappa_i,\kappa_{\mathrm{can}},
                         \kappa_{\mathcal L}\}
\]
proves the asserted two-outcome alternative on every scale.
\end{proof}

\begin{theorem}[Bounded curvature at bounded distance]
\label{thm:bounded-curvature-at-bounded-distance}
\dcref{MT:ch10:10.2}
\notready
For sufficiently small $\epsilon>0$ and every $A,C<\infty$, there are
$D_0=D_0(A,C,\epsilon)$ and $D=D(A,C,\epsilon)$ with the following
property.  Let $x$ lie in a three-dimensional Ricci flow with surgery whose
curvature is pinched toward positive.  Suppose every point $z$ at time at
most $\mathbf t(x)$ with $R(z)\ge4R(x)$ has a strong
$(C,\epsilon)$-canonical neighborhood.  If $R(x)\ge D_0$, then
\[
 R(y)\le D R(x)
 \qquad\text{for every }y\in
 B_{g(\mathbf t(x))}\bigl(x,A R(x)^{-1/2}\bigr).
\]
\end{theorem}

\begin{proof}
\uses{def:ricci-flow-with-surgery,def:strong-canonical-neighborhood-control,def:epsilon-neck-and-scale,lem:epsilon-neck-metric-control,prop:overlapping-epsilon-necks,lem:neck-chain-product,lem:nonflat-terminal-cone-exclusion}
Rescale so that $R(x)=1$ and suppose that points at bounded distance have
unbounded scalar curvature.  Along a minimizing geodesic from $x$, choose
the first points at which successive curvature levels are reached.  Above
level four, the canonical-neighborhood hypothesis applies.  A cap or compact
component at one of these first points has controlled diameter and curvature
ratio, so it cannot contain both the preceding controlled point and a point
at the next arbitrarily larger level.  The intervening regions are therefore
$\epsilon$-necks.

The metric and overlap estimates for necks assemble the relevant part of the
geodesic into a balanced chain.  By \ref{lem:neck-chain-product}, this chain
is a product tube.  Rescale at points whose neck radii tend to zero.  The
axial distance estimates and curvature pinching give a pointed incomplete
limit with nonnegative curvature, bounded curvature at one end, and
unbounded curvature at the other.  Blowing down the unbounded end, triangle
comparison shows that distances converge to
\[
 d_C((r,u),(s,v))^2=r^2+s^2-2rs\cos d_\Sigma(u,v),
\]
where the space of directions $\Sigma$ contains at least two points; the
limit is a genuine metric cone and not a ray.

The neck curvature derivative estimates provide uniform backward parabolic
neighborhoods on every compact subset away from the cone vertex.  Smooth
compactness therefore realizes the punctured cone as an open subset of a
nonnegatively curved Ricci flow at a positive time.  By
\ref{lem:nonflat-terminal-cone-exclusion} it is flat.  A flat punctured
three-cone has Euclidean volume ratios, so its vertex is removable and the
original curvature ratios remain bounded near the selected points.  This
contradicts their construction.  The contradiction gives a uniform $D$;
choosing $D_0$ large enough to make the pinching and neck estimates uniform
proves the theorem.
\end{proof}

\begin{theorem}[Pointed compactness for controlled flows]
\label{thm:pointed-controlled-flow-compactness}
\dcref{MT:ch11:11.1,MT:ch11:11.8,MT:ch11:11.10}
\notready
Let $(\mathcal M_j,G_j,x_j)$ be based controlled three-dimensional Ricci
flows and put $Q_j=R(x_j)\to\infty$.  Rescale by $Q_j$ and move the base
time to zero.  Assume curvature is pinched toward positive, the rescaled
flows are uniformly $\kappa$-noncollapsed on every fixed scale, and every
point at nonpositive time with rescaled scalar curvature at least four has a
strong $(C,\epsilon)$-canonical neighborhood.  Suppose that for some
$0<T_0\le\infty$, every fixed
$B(x_j,0,A)\times(-T,0]$, with $A<\infty$ and $T<T_0$, eventually has a
compatible spacetime embedding.  Then a subsequence converges smoothly to a
complete pointed Ricci flow on $(-T_0,0]$ with nonnegative curvature,
locally bounded in time and bounded on each time slice.  If $T_0=\infty$,
the limit is a $\kappa$-solution.
\end{theorem}

\begin{proof}
\uses{def:surgery-flow-noncollapsing,def:controlled-ricci-flow-with-surgery,thm:hamilton-compactness-generalized-flows,thm:bounded-curvature-at-bounded-distance}
On each fixed zero-time ball,
\ref{thm:bounded-curvature-at-bounded-distance} gives a scalar-curvature
bound.  Pinching converts it to a full curvature bound, and the canonical
derivative estimates propagate the bound for a uniform backward time.
The noncollapsing bound gives the basepoint volume hypothesis in
\ref{thm:hamilton-compactness-generalized-flows}.  Diagonalizing over $A$
and compact backward-time intervals gives a smooth limit complete at time
zero.
The negative part of the curvature operator vanishes after division by
$Q_j$, hence the limit has nonnegative curvature; the final clause of the
compactness theorem now makes every earlier slice complete.

To continue the limit backward, repeat this argument at the earliest time of
the current limit using neighborhoods already present in the approximating
flows.  If curvature is bounded there, Hamilton compactness extends the
limit.  If it is unbounded, choose first blow-up points.  The
bounded-curvature-at-bounded-distance theorem and the canonical-neighborhood
derivative estimates again give a uniform parabolic neighborhood, contrary
to maximality of the chosen interval.  This reaches every $T<T_0$.
The same bounded-distance argument rules out unbounded curvature on a fixed
limit slice.  When $T_0=\infty$, completeness, nonflatness,
nonnegative bounded curvature, and the assumed uniform noncollapsing are
precisely the defining properties of a $\kappa$-solution.
\end{proof}
